% ------------------------------------------------------------------------------
% Chapter 2
% ------------------------------------------------------------------------------
\chapter{Assertions} 
\label{cha:assertions}

In both code listing the clocking is excluded form the properties using a clocking block.
The binding uses all the signals of the dut as inputs and is connected to the module itself instead of the dut.
The fetch unit required to have access to internal logic, thi shas been enabled by the use of the prefix "\texttt{fetch\_unit.}".

\section{fetch\_unit}

Here we show the assertion used to test the fetch stage. 
In particular there are 6 properties each time dependant apart from the pc signal. 
The test intent is described in table \ref{tab:assertions_fetch}.
It can be observed that \texttt{P04} is not disabled by the reset as it is not syncronized (combinational).

\begin{table}[h]
	\centering
	\begin{tabular}{|l|l|l|}
		\hline
		\textbf{\#} & \textbf{Assertion} & \textbf{Intent} \\ \hline
		P00 		& Reset behavior 	 & On reset, PC must go to 0 \\ \hline
		P01 		& Branch priority	 & If \texttt{branch\_en}, PC must load \texttt{branch\_addr} regardless of \texttt{pc\_en} \\ \hline
		P02 		& PC increment   	 & If \texttt{pc\_en} is asserted (and no branch), PC must increment by 2 \\ \hline
		P03 		& PC stability   	 & If neither \texttt{branch\_en} nor \texttt{pc\_en} is asserted, PC must hold its value \\ \hline
		P04 		& Instr consistency  & \texttt{instr} must always equal mem[PC] \\ \hline
		P05 		& PC range safety 	 & PC must never go out of memory bounds (0-255) \\ \hline
	\end{tabular}
	\caption{Assertions table for the fetch stage}
	\label{tab:assertions_fetch}
\end{table}

\newpage
\begin{verbatim}
	property P00;
		!rst_n |=> fetch_unit.pc == 0;
	endproperty

	property P01;
		disable iff (!rst_n)
		branch_en |=> fetch_unit.pc == $past(branch_addr);
	endproperty

	property P02;
		disable iff (!rst_n)
		(pc_en && !branch_en) |=> fetch_unit.pc == $past(fetch_unit.pc) + 2;
	endproperty

	property P03;
		disable iff (!rst_n)
		(!pc_en && !branch_en) |=> fetch_unit.pc == $past(fetch_unit.pc);
	endproperty

	property P04;
		instr == fetch_unit.mem[fetch_unit.pc];
	endproperty

	property P05;
		disable iff (!rst_n)
		!$isunknown(fetch_unit.pc);
	endproperty
\end{verbatim}

\section{decode\_unit}

In the decode unit the first assertion developed followed the specification resulting in many errors. 
Following the dut implementationit was possible to find th ecorrect assertions.
The test intent is described in table \ref{tab:assertions_decode}.

\begin{table}[h]
	\centering
	\begin{tabular}{|l|l|l|}
		\hline
		\textbf{\#} & \textbf{Assertion} 		& \textbf{Intent} \\ \hline
		P00 		& Reset behavior 	 		& On reset, all outputs must got to 0 \\ \hline
		P01			& Accepted instruction		& On acceptance without hazard, all outputs should be set \\ \hline
		P02         & Stall						& On stall, all outputs should be stable \\ \hline
		P03			& Hazard					& On hazard, the flag should be high \\ \hline
		P04			& Decode delay (fwd)		& On acceptance, decode flag should be high after 2 cycles \\ \hline
		P05 		& Decode delay (bkw)		& Decode flag should be a pulse and enabled only on acceptance \\ \hline
		P06			& Decode (hazard)			& On hazard, non decode flag shpould be high \\ \hline
	\end{tabular}
	\caption{Assertions table for the decode stage}
	\label{tab:assertions_decode}
\end{table}

\begin{verbatim}
	property P00;
        !rst_n |=> (
            opcode       == 0 && 
            rd           == 0 && 
            rs           == 0 && 
            imm          == 0 && 
            decode_done  == 0 && 
            hazard_stall == 0
        );
    endproperty

    property P01;
        disable iff (!rst_n)
        instr_valid |=> 
          !hazard_stall |-> (
              opcode == $past(instr[15 : 12]) && 
              rd     == $past(instr[11 :  8]) && 
              rs     == $past(instr[ 7 :  4]) && 
              imm    == $past(instr[ 3 :  0])
          );
    endproperty


	
    property P02;
        disable iff (!rst_n)
        hazard_stall |-> (
            $stable(opcode) && 
            $stable(rd)     && 
            $stable(rs)     && 
            $stable(imm)
        );
    endproperty
    
    property P03;
        disable iff (!rst_n)
        (instr_valid && (instr[7:4] == rd)) |=> hazard_stall;
    endproperty

    property P04;
        disable iff (!rst_n)
        instr_valid |=> !hazard_stall |-> ##2 decode_done;
    endproperty

    property P05;
        disable iff (!rst_n)
        decode_done |-> $past(!hazard_stall, 2) && $past(instr_valid, 3);
    endproperty

    property P06;
        disable iff (!rst_n)
        hazard_stall |-> !decode_done;
    endproperty
\end{verbatim}